\documentclass[11pt,a4paper]{article}

\usepackage[brazil]{babel}
\usepackage[utf8]{inputenc}
\usepackage[T1]{fontenc}
\usepackage{lmodern}
\usepackage{geometry}
\usepackage{setspace}
\usepackage{parskip}
\usepackage{titlesec}
\usepackage{xcolor}
\usepackage{hyperref}

\geometry{margin=2.5cm}
\setstretch{1.2}

\definecolor{primary}{RGB}{41,128,185}
\hypersetup{
    colorlinks=true,
    linkcolor=primary,
    urlcolor=primary
}

\titleformat{\section}{\Large\bfseries\color{primary}}{\thesection}{1em}{}
\titleformat{\subsection}{\large\bfseries}{\thesubsection}{1em}{}

\title{\textbf{Composição da Chapa 01 - "stats.sh" \\ GULECD/UFPR}}
\author{Chapa 01}
\date{\today}

\begin{document}

\maketitle

\section{Coordenação Representativa}

\subsection{Coordenador Representativo}

\textbf{Nome:} Marcos Oliveira de Carvalho

\textbf{Curso/Vínculo:} Estatística e Ciência de Dados - UFPR

\textbf{Mini-biografia:}  
Marcos de Carvalho é farmacêutico (UFSM), com mestrado em Biologia Celular e Molecular (UFRGS), doutorado em Genética e Biologia Molecular (UFRGS) e formação executiva pelo Babson College (MA/EUA). Possui experiência em bioinformática e biologia computacional, com atuação em áreas como genômica, metagenômica e oncogenômica. Foi fundador e líder de duas startups de biotecnologia, com trajetória voltada à inovação e desenvolvimento de soluções baseadas em dados integradas a processos biotecnológicos. Atualmente atua como pesquisador de pós-doutorado na área de medicina translacional no CHC-UFPR, com interesse em integração de dados clínicos e moleculares, ao mesmo tempo em que cursa graduação em Estatística e Ciência de Dados (UFPR). Tem grande interesse em software livre/código aberto e ciência/tecnologia. É usuário de sistemas Linux desde 1999, com experiência em computadores pessoais e servidores até clusters de computação de alta performance. Apaixonado por música e metalhead desde criança.

\vspace{1em}

\subsection{Vice-Coordenador Representativo}

\textbf{Nome:} Rafael Alisson Schipiura

\textbf{Curso/Vínculo:} Estatística e Ciência de Dados - UFPR

\textbf{Mini-biografia:}  
Rafael Alisson Schipiura é Técnico em Engenharia de Dados na Celepar, com interesse em engenharia de dados, infraestrutura e soluções open source. Possui experiência de 13 anos como Administrador de Sistemas Desktop Linux, além de atuação recente em análise de dados, pipelines de dados e suporte a ambientes corporativos Linux. Atualmente cursa sua primeira graduação.

\section{Núcleo de Coordenação - Em ordem alfabética}

\subsection{Membro do Núcleo de Coordenação 1}

\textbf{Nome:} Gabriel Borges

\textbf{Cargo Proposto:} Coordenador Associado

\textbf{Curso/Vínculo:} Estatística e Ciência de Dados - UFPR

\textbf{Mini-biografia:}  
Gabriel Borges é Cientista de Dados, com forte interesse em aprendizado de máquina clássico, agentes conversacionais, Reinforcement Learning, Federated Learning, detecção de anomalias e visão computacional. Atualmente na VTEX, possui experiência e atuação no desenvolvimento e no ciclo de vida de modelos de ML em produção, observabilidade e escalabilidade de sistemas de IA, pesquisa aplicada em saúde digital, liderança da equipe de Ciência de Dados em projeto de extensão tecnológica financiado pela FACEPE, além de diversas publicações científicas. Também possui experiência em programação em Python, C, Java e Typescript, manutenção de biblioteca publicada no PyPI, contribuição para a comunidade Arch Linux por meio do AUR, uso avançado de sistemas Linux, análise de dados e participação em pesquisa em Visão Computacional e Inteligência Artificial. É discente do curso de Estatística e Ciência de Dados na UFPR e membro estudante do IEEE.

\vspace{1em}

\subsection{Membro do Núcleo de Coordenação 2}

\textbf{Nome:} Guilherme dos Santos Verissimo

\textbf{Cargo Proposto:} Coordenador Associado

\textbf{Curso/Vínculo:} Estatística e Ciência de Dados - UFPR

\textbf{Mini-biografia:}  
Guilherme dos Santos Verissimo é estudante de Estatística e Ciência de Dados da UFPR e atua profissionalmente há mais de três anos na área de tecnologia, com foco em engenharia de dados, automação e integração de sistemas corporativos. Foi co-fundador de uma empresa de marketing onde coordenou o time de tecnologia e liderou o desenvolvimento de sistemas internos. Atualmente se especializa em Ciência e Engenharia de Dados com o objetivo de viabilizar o acesso à informação de qualidade para públicos de diferentes áreas e níveis de formação. Em todos os seus projetos utiliza Linux como sistema operacional e prioriza o uso de software livre, por acreditar na urgência da democratização do acesso à tecnologia e na necessidade de soberania tecnológica do Brasil frente às grandes corporações do setor.
\vspace{1em}

\subsection{Membro do Núcleo de Coordenação 3}

\textbf{Nome:} Matheus Cabús

\textbf{Cargo Proposto:} Coordenador Associado

\textbf{Curso/Vínculo:} Estatística e Ciência de Dados - UFPR

\textbf{Mini-biografia:}  
Matheus Cabús é estudante de Estatística e Ciência de Dados pela UFPR com forte curiosidade interdisciplinar. Domina R, Python e SQL para análise, modelagem estatística e machine learning, tem experiência em programação (Python, C/C++), desenvolvimento web (HTML, CSS, JavaScript), além de conhecimentos em eletrônica básica (Arduino, PCBs) e paixão por produção de conteúdo como fotografia, edição de imagem e vídeo, design vetorial e produção musical. Atualmente está trabalhando como Estagiário na Solution IPD e Consultor de Informática na QCP.
\vspace{1em}

\subsection{Membro do Núcleo de Coordenação 4}

\textbf{Nome:} Rafael Noronha

\textbf{Cargo Proposto:} Coordenador Associado

\textbf{Curso/Vínculo:} Estatística e Ciência de Dados - UFPR

\textbf{Mini-biografia:}  
Rafael Noronha é desenvolvedor de software com atuação em backend. Trabalha com sistemas em cloud, utilizando principalmente Golang, PHP e Node.js Tem interesse em infraestrutura e engenharia de confiabilidade em software (SRE), além de atuar com ambientes Linux no dia a dia, possui experiência com desenvolvimento de APIs, integração entre serviços e uso de bancos de dados relacionais.

\vspace{1em}

\end{document}